\documentclass[10pt,a4paper]{article}
\usepackage{../../styles/neet_style}

\pagestyle{fancy}
\fancyhf{}
\fancyhead[L]{\small \textbf{NEET UG 2027: Thermal Physics \& Oscillations Master Theory}}
\fancyhead[R]{\small Satya Prakash}
\fancyfoot[L]{\small \textbf{NEET UG 2027 Prep Series}}
\fancyfoot[C]{\thepage}
\fancyfoot[R]{\small \textbf{Physics Chapter Notes}}

\begin{document}

% Title Header Banner
\begin{center}
  {\Huge \textbf{NEET UG 2027: Master Chapter Theory}}\\[6pt]
  {\Large \textbf{Thermal Physics, Kinetic Theory, Thermodynamics \& Oscillations}}\\[4pt]
  {\small \textbf{Prepared, Compiled \& Authored by Satya Prakash} $\cdot$ Strictly Calibrated to Revised NMC Syllabus}
  \vspace{4pt}
  \hrule height 1.5pt
\end{center}

\section{Thermal Properties of Matter, Thermal Expansion \& Calorimetry}

\begin{conceptbox}[Foundations of Heat, Temperature \& Thermometry]
\textbf{1. Heat vs. Temperature:}
\begin{itemize}
  \item \textbf{Heat ($Q$):} Energy in transit across a boundary solely driven by a macroscopic temperature gradient. Heat is a path-dependent boundary phenomenon, not an intrinsic property stored within a system ($[Q] = \mathrm{J}$ or $\mathrm{cal}$, where $1\,\mathrm{cal} = 4.186\,\mathrm{J}$).
  \item \textbf{Temperature ($T$):} Macroscopic measure of the average translational kinetic energy per molecule of a substance. It dictates the direction of spontaneous thermal energy transfer (from higher $T$ to lower $T$).
\end{itemize}
\textbf{2. Thermometric Scales \& Interconversions:}
Any linear thermometric property $X(T)$ (e.g., liquid column length, platinum wire resistance, gas pressure at constant volume) satisfies:
\begin{equation}
  \frac{T - \mathrm{LFP}}{\mathrm{UFP} - \mathrm{LFP}} = \text{constant} \implies \frac{C - 0}{100} = \frac{F - 32}{180} = \frac{K - 273.15}{100} = \frac{R - 0}{80}
\end{equation}
where $\mathrm{LFP}$ is the Lower Fixed Point (ice point at $1\,\mathrm{atm}$) and $\mathrm{UFP}$ is the Upper Fixed Point (steam point at $1\,\mathrm{atm}$).
\begin{equation}
  \frac{C}{5} = \frac{F - 32}{9} = \frac{K - 273.15}{5}
\end{equation}
\textit{Key Coincidence Points:}
\begin{itemize}
  \item $C = F = -40^\circ \implies -40^\circ\mathrm{C} = -40^\circ\mathrm{F}$.
  \item $F = K = 574.25 \implies 574.25^\circ\mathrm{F} = 574.25\,\mathrm{K}$.
  \item Incremental temperature step equivalence: $\Delta T = 1\,\mathrm{K} \equiv 1^\circ\mathrm{C} \equiv 1.8^\circ\mathrm{F}$.
\end{itemize}
\end{conceptbox}

\begin{formulabox}[Thermal Expansion in Solids \& Liquids: Master Governing Equations]
\textbf{1. Anisotropic and Isotropic Thermal Expansion in Solids:}
\begin{itemize}
  \item \textbf{Linear Expansion:}
  \begin{equation}
    L(T) = L_0 (1 + \alpha \Delta T) \implies \Delta L = L_0 \alpha \Delta T, \quad \text{where } \alpha = \frac{1}{L_0}\frac{dL}{dT} \quad [\mathrm{K}^{-1} \text{ or } ^\circ\mathrm{C}^{-1}]
  \end{equation}
  \item \textbf{Superficial (Areal) Expansion:}
  \begin{equation}
    A(T) = A_0 (1 + \beta \Delta T) \implies \Delta A = A_0 \beta \Delta T, \quad \text{where } \beta \approx 2\alpha
  \end{equation}
  \item \textbf{Volumetric (Cubical) Expansion:}
  \begin{equation}
    V(T) = V_0 (1 + \gamma \Delta T) \implies \Delta V = V_0 \gamma \Delta T, \quad \text{where } \gamma \approx 3\alpha
  \end{equation}
  \item \textbf{Isotropic Ratio:} $\alpha : \beta : \gamma = 1 : 2 : 3 \iff \alpha = \frac{\beta}{2} = \frac{\gamma}{3}$.
\end{itemize}
\textbf{2. Proof of $\gamma = 3\alpha$ for Isotropic Solids:}
Consider a cube of initial edge length $L_0$ and initial volume $V_0 = L_0^3$. When temperature rises by $\Delta T$:
\begin{equation}
  V = L^3 = [L_0(1 + \alpha\Delta T)]^3 = L_0^3(1 + 3\alpha\Delta T + 3\alpha^2\Delta T^2 + \alpha^3\Delta T^3)
\end{equation}
Since $\alpha \sim 10^{-5}\,\mathrm{K}^{-1}$, higher-order terms $\alpha^2\Delta T^2 \ll 1$ and $\alpha^3\Delta T^3 \ll 1$:
\begin{equation}
  V \approx V_0(1 + 3\alpha\Delta T) \equiv V_0(1 + \gamma\Delta T) \implies \mathbf{\gamma = 3\alpha}
\end{equation}
\textbf{3. Apparent vs. Real Expansion of Liquids in Vessels:}
\begin{equation}
  \gamma_{\text{real}} = \gamma_{\text{apparent}} + \gamma_{\text{vessel}}
\end{equation}
\begin{itemize}
  \item If $\gamma_{\text{real}} > \gamma_{\text{vessel}}$, liquid level rises upon heating.
  \item If $\gamma_{\text{real}} < \gamma_{\text{vessel}}$, liquid level falls upon heating.
  \item If $\gamma_{\text{real}} = \gamma_{\text{vessel}}$, liquid level remains invariant regardless of temperature.
\end{itemize}
\textbf{4. Density Variation with Temperature:}
Let mass $M$ occupy volume $V_0$ at $T_0$ ($\rho_0 = M/V_0$). At temperature $T_0 + \Delta T$:
\begin{equation}
  \rho(T) = \frac{M}{V(T)} = \frac{M}{V_0(1 + \gamma\Delta T)} = \rho_0(1 + \gamma\Delta T)^{-1} \approx \mathbf{\rho_0(1 - \gamma\Delta T)}
\end{equation}
\end{formulabox}

\begin{usecasebox}[Thermal Stress in Rigidly Clamped Rods \& Bimetallic Strips]
\textbf{1. Thermal Stress Formulation:}
When a rod of length $L_0$, cross-sectional area $A$, and Young's modulus $Y$ is clamped rigidly between two immovable walls and heated by $\Delta T$:
\begin{itemize}
  \item Free thermal expansion prevented: $\Delta L = L_0 \alpha \Delta T$.
  \item Compressive thermal strain induced: $\varepsilon_{\text{thermal}} = \frac{\Delta L}{L_0} = \alpha \Delta T$.
  \item Compressive stress developed:
  \begin{equation}
    \mathbf{\sigma_{\text{th}} = \frac{F}{A} = Y \varepsilon_{\text{thermal}} = Y \alpha \Delta T}
  \end{equation}
  \item Thermal force exerted on clamps: $\mathbf{F = Y A \alpha \Delta T}$. Note that $\sigma_{\text{th}}$ is completely independent of the rod length $L_0$.
\end{itemize}
\textbf{2. Bimetallic Strip Curvature:}
Two strips of equal thickness $d$, made of metals with expansion coefficients $\alpha_1 > \alpha_2$, are riveted together. Upon heating, metal 1 expands more than metal 2, forcing the composite strip into an arc of radius $R$:
\begin{equation}
  R \approx \frac{d}{(\alpha_1 - \alpha_2)\Delta T}
\end{equation}
The strip bends with the higher-$\alpha$ metal on the convex outer side. Upon cooling ($\Delta T < 0$), it bends in the reverse direction (higher-$\alpha$ metal on the concave inner side).
\end{usecasebox}

\begin{formulabox}[Calorimetry, Specific Heat Capacity \& Latent Heat]
\textbf{1. Heat Capacities:}
\begin{itemize}
  \item \textbf{Specific Heat Capacity ($s$ or $c$):} Heat required per unit mass to raise temperature by $1\,\mathrm{K}$:
  \begin{equation}
    dQ = m s dT \implies Q = m \int_{T_1}^{T_2} s(T) dT = m s \Delta T \quad (\text{if } s = \text{const})
  \end{equation}
  For water: $s_w = 1.0\,\mathrm{cal}/(\mathrm{g}\cdot^\circ\mathrm{C}) = 4186\,\mathrm{J}/(\mathrm{kg}\cdot\mathrm{K})$. For ice: $s_{\text{ice}} \approx 0.5\,\mathrm{cal}/(\mathrm{g}\cdot^\circ\mathrm{C}) = 2100\,\mathrm{J}/(\mathrm{kg}\cdot\mathrm{K})$.
  \item \textbf{Molar Heat Capacity ($C$):} $C = \frac{1}{n}\frac{dQ}{dT} = M_{\text{molar}} s$.
  \item \textbf{Water Equivalent of Calorimeter ($W$):} Mass of water that absorbs/releases the same heat as the calorimeter for an identical temperature change:
  \begin{equation}
    W = m_{\text{cal}} \frac{s_{\text{cal}}}{s_w} \implies Q = (m_w + W) s_w \Delta T
  \end{equation}
\end{itemize}
\textbf{2. Latent Heat \& Phase Changes (Isothermal Heat Addition):}
\begin{equation}
  Q = m L
\end{equation}
\begin{itemize}
  \item \textbf{Latent Heat of Fusion of Ice ($L_f$):} $L_f \approx 80\,\mathrm{cal}/\mathrm{g} = 3.33 \times 10^5\,\mathrm{J}/\mathrm{kg}$ at $0^\circ\mathrm{C}$.
  \item \textbf{Latent Heat of Vaporization of Water ($L_v$):} $L_v \approx 540\,\mathrm{cal}/\mathrm{g} = 2.26 \times 10^6\,\mathrm{J}/\mathrm{kg}$ at $100^\circ\mathrm{C}$.
\end{itemize}
\textbf{3. Principle of Calorimetry (Conservation of Energy):}
In a thermally isolated calorimeter:
\begin{equation}
  \mathbf{\sum Q_{\text{lost by hot bodies}} = \sum Q_{\text{gained by cold bodies}}}
\end{equation}
\end{formulabox}

\begin{trapbox}[Calorimetry Pitfalls: Ice-Water-Steam Equilibrium Traps]
\textbf{Common Error:} Blindly applying $m_1 s_1 (T_1 - T_f) = m_2 s_2 (T_f - T_2)$ without verifying whether sufficient heat exists to melt all ice or condense all steam.
\begin{itemize}
  \item \textbf{Ice + Water Scenario:} Always calculate $Q_{\text{available}}$ from water cooling to $0^\circ\mathrm{C}$ vs. $Q_{\text{melt}}$ required to melt all ice at $0^\circ\mathrm{C}$. If $Q_{\text{available}} < m_{\text{ice}} L_f$, the final mixture temperature is \textbf{strictly $0^\circ\mathrm{C}$}, and only a fraction $m' = Q_{\text{available}}/L_f$ of ice melts.
  \item \textbf{Steam + Water Scenario:} Always verify if water can absorb all latent heat of steam ($540\,\mathrm{cal}/\mathrm{g}$) before heating to $100^\circ\mathrm{C}$. If not, final equilibrium temperature remains \textbf{strictly $100^\circ\mathrm{C}$}.
\end{itemize}
\end{trapbox}

\section{Modes of Heat Transfer \& Radiation Laws}

\begin{conceptbox}[Thermal Conduction: Fourier's Law \& Composite Slabs]
\textbf{1. Fourier's Law of Steady-State Conduction:}
Thermal energy flowing per unit time through a cross-sectional area $A$ normal to a temperature gradient $dT/dx$:
\begin{equation}
  \mathbf{H = \frac{dQ}{dt} = -k A \frac{dT}{dx}}
\end{equation}
where $k$ is thermal conductivity ($\mathrm{W}\cdot\mathrm{m}^{-1}\cdot\mathrm{K}^{-1}$). For a uniform rod of length $L$ with ends maintained at $T_H$ and $T_C$:
\begin{equation}
  H = \frac{k A (T_H - T_C)}{L} = \frac{T_H - T_C}{R_{\text{th}}}, \quad \text{where } \mathbf{R_{\text{th}} = \frac{L}{k A}} \quad [\mathrm{K}\cdot\mathrm{W}^{-1}]
\end{equation}
$R_{\text{th}}$ is the \textbf{thermal resistance}, directly analogous to electrical resistance $R_{\text{elec}} = \frac{L}{\sigma A}$.
\end{conceptbox}

\begin{formulabox}[Series \& Parallel Thermal Conductors: Derivations]
\textbf{1. Series Combination of Two Conducting Slabs:}
Consider two slabs of lengths $L_1, L_2$, areas $A$, and conductivities $k_1, k_2$ placed in end-to-end contact between temperatures $T_1$ and $T_2$ ($T_1 > T_2$). Let $T_0$ be the interface temperature. Under steady state, the heat current $H$ is identical through both:
\begin{equation}
  H = \frac{k_1 A (T_1 - T_0)}{L_1} = \frac{k_2 A (T_0 - T_2)}{L_2} \implies R_{\text{eq}} = R_1 + R_2 = \frac{L_1}{k_1 A} + \frac{L_2}{k_2 A}
\end{equation}
Equating expressions to solve for the interface temperature $T_0$:
\begin{equation}
  \mathbf{T_0 = \frac{\frac{k_1}{L_1} T_1 + \frac{k_2}{L_2} T_2}{\frac{k_1}{L_1} + \frac{k_2}{L_2}} = \frac{k_1 L_2 T_1 + k_2 L_1 T_2}{k_1 L_2 + k_2 L_1}}
\end{equation}
Equivalent thermal conductivity $k_{\text{eq}}$ for total length $L_1 + L_2$:
\begin{equation}
  \frac{L_1 + L_2}{k_{\text{eq}} A} = \frac{L_1}{k_1 A} + \frac{L_2}{k_2 A} \implies \mathbf{k_{\text{eq}} = \frac{L_1 + L_2}{\frac{L_1}{k_1} + \frac{L_2}{k_2}}} \quad \left(\text{If } L_1 = L_2: \; k_{\text{eq}} = \frac{2k_1 k_2}{k_1 + k_2}\right)
\end{equation}
\textbf{2. Parallel Combination of Two Conducting Slabs:}
Plates of areas $A_1, A_2$ and common length $L$ connected between reservoirs at $T_1$ and $T_2$:
\begin{equation}
  H_{\text{net}} = H_1 + H_2 \implies \frac{1}{R_{\text{eq}}} = \frac{1}{R_1} + \frac{1}{R_2} \implies \mathbf{k_{\text{eq}} = \frac{k_1 A_1 + k_2 A_2}{A_1 + A_2}}
\end{equation}
\end{formulabox}

\begin{conceptbox}[Thermal Radiation: Stefan-Boltzmann \& Wien's Laws]
\textbf{1. Blackbody Radiation \& Prevost's Theory:}
Prevost's theory asserts that all bodies at temperatures $T > 0\,\mathrm{K}$ emit thermal radiation continuously to, and absorb radiation from, their surroundings. Thermal equilibrium represents dynamic parity between emission and absorption.
\begin{itemize}
  \item \textbf{Stefan-Boltzmann Law:} Total radiant energy emitted per unit surface area per second by a blackbody:
  \begin{equation}
    E_b = \sigma T^4
  \end{equation}
  where $\sigma = 5.670 \times 10^{-8}\,\mathrm{W}\cdot\mathrm{m}^{-2}\cdot\mathrm{K}^{-4}$ is the Stefan-Boltzmann constant.
  \item For a non-blackbody of surface area $A$ and emissivity $e$ ($0 < e < 1$):
  \begin{equation}
    P_{\text{emit}} = e \sigma A T^4
  \end{equation}
  \item \textbf{Net Radiative Heat Transfer with Surroundings at $T_0$:}
  \begin{equation}
    \mathbf{P_{\text{net}} = \frac{dQ}{dt} = e \sigma A (T^4 - T_0^4)}
  \end{equation}
  \item \textbf{Wien's Displacement Law:} The wavelength $\lambda_m$ corresponding to maximum spectral emissive power is inversely proportional to absolute temperature:
  \begin{equation}
    \mathbf{\lambda_m T = b = 2.898 \times 10^{-3}\,\mathrm{m}\cdot\mathrm{K}}
  \end{equation}
  As temperature increases, peak emission shifts towards shorter wavelengths (higher frequencies).
\end{itemize}
\end{conceptbox}

\begin{formulabox}[Newton's Law of Cooling: Derivation \& Analytical Form]
\textbf{1. Calculus Derivation from Stefan-Boltzmann Law:}
Consider a body of mass $m$, specific heat capacity $s$, and surface area $A$ at temperature $T$ cooling in surroundings at temperature $T_0$, where $\Delta T = T - T_0 \ll T_0$.
\begin{equation}
  T^4 - T_0^4 = (T_0 + \Delta T)^4 - T_0^4 = T_0^4 \left[\left(1 + \frac{\Delta T}{T_0}\right)^4 - 1\right]
\end{equation}
Using the binomial expansion $(1 + x)^4 \approx 1 + 4x$ for $x = \Delta T / T_0 \ll 1$:
\begin{equation}
  T^4 - T_0^4 \approx T_0^4 \left(1 + 4\frac{\Delta T}{T_0} - 1\right) = 4 T_0^3 (T - T_0)
\end{equation}
Substituting this into the net radiant heat loss rate:
\begin{equation}
  \frac{dQ}{dt} = -m s \frac{dT}{dt} = e \sigma A \cdot 4 T_0^3 (T - T_0) \implies \mathbf{\frac{dT}{dt} = -k (T - T_0)}
\end{equation}
where cooling constant $k = \frac{4 e \sigma A T_0^3}{m s}$.
\textbf{2. Exact Exponential Solution:}
\begin{equation}
  \int_{T_i}^{T(t)} \frac{dT}{T - T_0} = -k \int_0^t dt \implies \mathbf{T(t) - T_0 = (T_i - T_0) e^{-k t}}
\end{equation}
\textbf{3. Practical Average Approximation Formula for Problems:}
If a body cools from $T_1$ to $T_2$ in time $t$:
\begin{equation}
  \mathbf{\frac{T_1 - T_2}{t} = K \left(\frac{T_1 + T_2}{2} - T_0\right)}
\end{equation}
\end{formulabox}

\begin{ncertbox}[NCERT Extraction: Anomalous Expansion of Water]
Water exhibits anomalous expansion between $0^\circ\mathrm{C}$ and $4^\circ\mathrm{C}$. When water is heated from $0^\circ\mathrm{C}$, its volume \textbf{decreases} and its density \textbf{increases}, reaching a maximum density of $\rho_{\max} \approx 1.000\,\mathrm{g/cm}^3$ at $3.98^\circ\mathrm{C} \approx 4^\circ\mathrm{C}$. Above $4^\circ\mathrm{C}$, water expands normally with temperature.
\begin{itemize}
  \item \textit{Ecological Consequence:} In freezing winters, surface water cools to $4^\circ\mathrm{C}$ and sinks. Once the entire body reaches $4^\circ\mathrm{C}$, further cooling drops surface water to $0^\circ\mathrm{C}$ where it freezes into ice ($\rho_{\text{ice}} < \rho_{\text{water}}$). The ice layer floats on top, acting as a thermal insulator ($k_{\text{ice}} \approx 2.2\,\mathrm{W/m\cdot K}$), preserving liquid water at $4^\circ\mathrm{C}$ beneath and sustaining aquatic life.
\end{itemize}
\end{ncertbox}

\section{Thermodynamics: First Law, Processes \& Mayer's Relation}

\begin{conceptbox}[Thermodynamic Systems, State Functions \& First Law]
\textbf{1. Thermodynamic Equilibrium \& Zeroth Law:}
Two systems in mutual thermal equilibrium with a third system are in thermal equilibrium with each other. This establishes temperature as the unique scalar state property governing thermal equilibrium.
\textbf{2. Internal Energy ($U$):}
Internal energy is a unique thermodynamic \textbf{state function} reflecting the sum of molecular kinetic and intermolecular potential energies. For an ideal gas (zero intermolecular potential):
\begin{equation}
  U = U(T) = n C_v T \implies \mathbf{dU = n C_v dT}
\end{equation}
Across any cyclic process, $\Delta U_{\text{cycle}} = \oint dU = 0$.
\textbf{3. First Law of Thermodynamics (Energy Conservation):}
\begin{equation}
  \mathbf{\Delta Q = \Delta U + W} \iff \mathbf{dQ = dU + dW} = n C_v dT + P dV
\end{equation}
\textit{Physics Sign Conventions:}
\begin{itemize}
  \item $Q > 0$ when heat is absorbed by the system; $Q < 0$ when heat is released.
  \item $W = \int P dV > 0$ when work is done \textbf{by} the system (expansion, $dV > 0$); $W < 0$ when work is done \textbf{on} the system (compression, $dV < 0$).
\end{itemize}
\end{conceptbox}

\begin{formulabox}[Work Done in Master Thermodynamic Processes: Step-by-Step Derivations]
\textbf{1. Isochoric Process (Constant Volume, $V = \text{const}$):}
\begin{equation}
  dV = 0 \implies \mathbf{W = 0}, \quad \mathbf{\Delta Q = \Delta U = n C_v \Delta T}
\end{equation}
\textbf{2. Isobaric Process (Constant Pressure, $P = \text{const}$):}
\begin{equation}
  \mathbf{W = P \Delta V = P(V_2 - V_1) = n R (T_2 - T_1)}
\end{equation}
\begin{equation}
  \Delta Q = n C_p \Delta T, \quad \Delta U = n C_v \Delta T \implies \mathbf{\Delta U : W : \Delta Q = C_v : R : C_p = 1 : (\gamma - 1) : \gamma}
\end{equation}
\textbf{3. Isothermal Process (Constant Temperature, $T = \text{const}$):}
Since $T = \text{const}$, $\Delta U = 0 \implies \Delta Q = W$. Using $P = \frac{nRT}{V}$:
\begin{equation}
  \mathbf{W = \int_{V_1}^{V_2} P dV = n R T \int_{V_1}^{V_2} \frac{dV}{V} = n R T \ln\left(\frac{V_2}{V_1}\right) = n R T \ln\left(\frac{P_1}{P_2}\right) = 2.303 n R T \log_{10}\left(\frac{V_2}{V_1}\right)}
\end{equation}
\textbf{4. Adiabatic Process (Zero Heat Exchange, $\Delta Q = 0$):}
Governing equations: $P V^\gamma = \text{const}$, $T V^{\gamma-1} = \text{const}$, $P^{1-\gamma} T^\gamma = \text{const}$.
From First Law: $dQ = 0 \implies dW = -dU = -n C_v dT$.
\begin{equation}
  W = -\int_{T_1}^{T_2} n C_v dT = -n C_v (T_2 - T_1) = n \left(\frac{R}{\gamma - 1}\right)(T_1 - T_2)
\end{equation}
Using $nRT_1 = P_1 V_1$ and $nRT_2 = P_2 V_2$:
\begin{equation}
  \mathbf{W = \frac{P_1 V_1 - P_2 V_2}{\gamma - 1} = \frac{n R (T_1 - T_2)}{\gamma - 1} = -\Delta U}
\end{equation}
\textbf{5. General Polytropic Process ($P V^m = \text{const}$):}
\begin{equation}
  W = \frac{P_1 V_1 - P_2 V_2}{m - 1}, \qquad \mathbf{C = C_v + \frac{R}{1 - m} = \frac{R}{\gamma - 1} + \frac{R}{1 - m}}
\end{equation}
\end{formulabox}

\begin{formulabox}[Mayer's Relation Derivation: $C_p - C_v = R$]
Consider 1 mole of an ideal gas ($n=1$).
\begin{enumerate}
  \item Heat added at constant volume: $dQ_v = C_v dT$. Since $dV = 0$, $dW = 0$, so by First Law:
  \begin{equation}
    dU = dQ_v = C_v dT
  \end{equation}
  Because internal energy of an ideal gas depends solely on temperature, $dU = C_v dT$ holds across \textbf{any} process.
  \item Heat added at constant pressure: $dQ_p = C_p dT$. By First Law:
  \begin{equation}
    dQ_p = dU + dW = C_v dT + P dV
  \end{equation}
  \item Differentiating the ideal gas equation of state $P V = R T$ at constant pressure $P$:
  \begin{equation}
    P dV = R dT
  \end{equation}
  \item Substituting $P dV = R dT$ into the First Law expression:
  \begin{equation}
    C_p dT = C_v dT + R dT \implies (C_p - C_v) dT = R dT \implies \mathbf{C_p - C_v = R}
  \end{equation}
\end{enumerate}
\textit{In terms of specific heats (per unit mass):} $c_p - c_v = \frac{R}{M_{\text{molar}}}$.
\end{formulabox}

\begin{usecasebox}[P-V Indicator Diagram Slopes \& Cyclic Process Work]
\textbf{1. Comparison of Adiabatic vs. Isothermal Slopes:}
\begin{itemize}
  \item Isothermal curve: $P V = \text{const} \implies P dV + V dP = 0 \implies \mathbf{\left(\frac{dP}{dV}\right)_{\text{iso}} = -\frac{P}{V}}$.
  \item Adiabatic curve: $P V^\gamma = \text{const} \implies V^\gamma dP + \gamma P V^{\gamma-1} dV = 0 \implies \mathbf{\left(\frac{dP}{dV}\right)_{\text{ad}} = -\gamma \frac{P}{V}}$.
  \item \textbf{Fundamental Slope Theorem:}
  \begin{equation}
    \mathbf{\left(\frac{dP}{dV}\right)_{\text{adiabatic}} = \gamma \left(\frac{dP}{dV}\right)_{\text{isothermal}}}
  \end{equation}
  Since $\gamma > 1$, the adiabatic curve is \textbf{$\gamma$ times steeper} than the isothermal curve at any point of intersection.
\end{itemize}
\textbf{2. Expansion from Common Initial State $(P_0, V_0)$ to Final Volume $V_f > V_0$:}
\begin{equation}
  P_{\text{isobaric}} > P_{\text{isothermal}} > P_{\text{adiabatic}} \implies \mathbf{W_{\text{isobaric}} > W_{\text{isothermal}} > W_{\text{adiabatic}}}
\end{equation}
\textbf{3. Compression from Common Initial State $(P_0, V_0)$ to Final Volume $V_f < V_0$:}
Work done \textbf{on} the gas is largest for adiabatic compression:
\begin{equation}
  |W_{\text{adiabatic}}| > |W_{\text{isothermal}}| > |W_{\text{isobaric}}|
\end{equation}
\textbf{4. Net Work in Cyclic Processes:}
The net work in any closed cycle on a $P-V$ diagram equals the enclosed loop area:
\begin{equation}
  W_{\text{net}} = \oint P dV = \begin{cases} > 0 & \text{Clockwise cycle (heat engine, net work output)} \\ < 0 & \text{Counter-clockwise cycle (work absorbed)} \end{cases}
\end{equation}
\end{usecasebox}

\section{Second Law of Thermodynamics \& Syllabus Boundaries}

\begin{conceptbox}[Second Law of Thermodynamics: Clausius \& Kelvin-Planck Formulations]
The First Law dictates energy conservation but fails to predict the spontaneous direction of natural processes. The Second Law establishes irreversibility:
\begin{itemize}
  \item \textbf{Kelvin-Planck Statement:} It is impossible to construct a heat engine operating in a thermodynamic cycle that absorbs thermal energy from a single thermal reservoir and converts $100\%$ of this energy into mechanical work without rejecting any heat to a colder reservoir. (No engine can operate at $100\%$ thermal efficiency; perpetual motion machines of the second kind are impossible).
  \item \textbf{Clausius Statement:} It is impossible to construct a cyclic device that transfers heat spontaneously from a cooler body to a hotter body without requiring external mechanical work input.
  \item \textbf{Equivalence:} A violation of the Kelvin-Planck statement directly entails a violation of the Clausius statement, and vice versa.
\end{itemize}
\end{conceptbox}

\begin{trapbox}[Syllabus Boundary Alert: Carnot Engine \& Refrigerator Deletion]
\textbf{NMC REVISION WARNING FOR NEET UG:}
Under the updated NEET UG syllabus issued by the National Medical Commission (NMC):
\begin{itemize}
  \item \textbf{DELETED TOPICS:} The theoretical formulation of the \textbf{Carnot Engine}, Carnot cycle efficiency ($\eta = 1 - T_C/T_H$), Carnot theorem proofs, and \textbf{Refrigerators / Heat Pumps} (coefficient of performance $\beta = Q_C/W = T_C/(T_H - T_C)$) have been \textbf{formally removed}.
  \item \textbf{EXAM ORIENTATION:} Questions will \textbf{not} test Carnot cycle calculations. Candidates must focus strictly on General Thermodynamic Cycles (triangular, rectangular $P-V$ loops), First Law applications ($\Delta Q = \Delta U + W$), work in isothermal/adiabatic/isobaric expansions, and the conceptual statements of the Second Law.
\end{itemize}
\end{trapbox}

\section{Kinetic Theory of Gases (KTG): Molecular Pressure \& Equipartition}

\begin{conceptbox}[Postulates of the Kinetic Theory of an Ideal Gas]
An ideal gas is modeled micro-mechanically under the following foundational postulates:
\begin{enumerate}
  \item The gas comprises an enormous number of identical, tiny, hard spherical molecules in ceaseless, chaotic random motion across all directions with all possible speeds.
  \item The physical volume occupied by the gas molecules themselves is negligibly small compared to the total macroscopic volume of the container ($V_{\text{molecules}} \ll V_{\text{container}}$).
  \item No intermolecular forces (attractive or repulsive) act between molecules except during collisions. Hence, potential energy is strictly zero; all internal energy is purely translational/rotational kinetic energy.
  \item All intermolecular collisions and collisions with container walls are perfectly elastic ($\Delta K = 0$). Collision durations are negligible compared to transit times between collisions.
  \item Gravity has negligible effect on molecular trajectory between collisions due to high speeds ($\sim 10^2 - 10^3\,\mathrm{m/s}$).
\end{enumerate}
\end{conceptbox}

\begin{formulabox}[First-Principles Derivation of Ideal Gas Pressure: $P = \frac{1}{3}\rho v_{\text{rms}}^2$]
Consider $N$ identical gas molecules, each of mass $m$, enclosed in a cubical box of side $L$ and volume $V = L^3$.
\begin{enumerate}
  \item Let a molecule possess velocity vector $\mathbf{v} = v_x \hat{\mathbf{i}} + v_y \hat{\mathbf{j}} + v_z \hat{\mathbf{k}}$ with speed squared $v^2 = v_x^2 + v_y^2 + v_z^2$.
  \item When this molecule collides elastically with the wall perpendicular to the $x$-axis at $x = L$:
  \begin{equation}
    p_{x,\text{initial}} = m v_x, \qquad p_{x,\text{final}} = -m v_x \implies \Delta p_x = -2 m v_x
  \end{equation}
  Momentum imparted to the wall per collision is $+2 m v_x$.
  \item Time elapsed between consecutive collisions of this molecule with the same wall:
  \begin{equation}
    \Delta t = \frac{2 L}{v_x}
  \end{equation}
  \item Average force exerted by this single molecule on the wall:
  \begin{equation}
    F_1 = \frac{\Delta p_{\text{wall}}}{\Delta t} = \frac{2 m v_x}{2 L / v_x} = \frac{m v_x^2}{L}
  \end{equation}
  \item Total force exerted by all $N$ molecules on the wall of area $A = L^2$:
  \begin{equation}
    F_{\text{total}} = \sum_{i=1}^N \frac{m v_{xi}^2}{L} = \frac{m}{L} N \overline{v_x^2}
  \end{equation}
  \item Since molecular motion is isotropic, $\overline{v_x^2} = \overline{v_y^2} = \overline{v_z^2} = \frac{1}{3}\overline{v^2} = \frac{1}{3} v_{\text{rms}}^2$.
  \item Pressure $P = \frac{F_{\text{total}}}{A} = \frac{m N \overline{v_x^2}}{L \cdot L^2} = \frac{1}{3}\frac{N m}{V} v_{\text{rms}}^2$:
  \begin{equation}
    \mathbf{P = \frac{1}{3}\rho v_{\text{rms}}^2 = \frac{1}{3}\frac{M_{\text{total}}}{V} v_{\text{rms}}^2}
  \end{equation}
\end{enumerate}
\end{formulabox}

\begin{formulabox}[Kinetic Interpretation of Temperature \& Molecular Speed Spectrum]
\textbf{1. Average Kinetic Energy per Molecule:}
Rewriting the pressure equation as $P V = \frac{2}{3} \left(N \cdot \frac{1}{2}m v_{\text{rms}}^2\right) = \frac{2}{3} E_{\text{trans}}$.
Equating to the ideal gas law $P V = N k_B T$ (where $k_B = R/N_A = 1.381 \times 10^{-23}\,\mathrm{J/K}$):
\begin{equation}
  \frac{2}{3} E_{\text{trans}} = N k_B T \implies \mathbf{\overline{K}_{\text{trans}} = \frac{1}{2}m v_{\text{rms}}^2 = \frac{3}{2}k_B T}
\end{equation}
Average translational kinetic energy depends \textbf{solely on absolute temperature $T$}, invariant of molecular mass or gas species!
\textbf{2. Molecular Speed Hierarchy:}
\begin{itemize}
  \item \textbf{Root-Mean-Square Speed ($v_{\text{rms}}$):}
  \begin{equation}
    \mathbf{v_{\text{rms}} = \sqrt{\frac{3 k_B T}{m}} = \sqrt{\frac{3 R T}{M_{\text{molar}}}} = \sqrt{\frac{3 P}{\rho}}}
  \end{equation}
  \item \textbf{Average Speed ($v_{\text{avg}}$):}
  \begin{equation}
    \mathbf{v_{\text{avg}} = \sqrt{\frac{8 k_B T}{\pi m}} = \sqrt{\frac{8 R T}{\pi M_{\text{molar}}}}} \approx 0.921 v_{\text{rms}}
  \end{equation}
  \item \textbf{Most Probable Speed ($v_{\text{mp}}$):}
  \begin{equation}
    \mathbf{v_{\text{mp}} = \sqrt{\frac{2 k_B T}{m}} = \sqrt{\frac{2 R T}{M_{\text{molar}}}}} \approx 0.816 v_{\text{rms}}
  \end{equation}
  \item \textbf{Golden Speed Ratio:}
  \begin{equation}
    \mathbf{v_{\text{mp}} : v_{\text{avg}} : v_{\text{rms}} = \sqrt{2} : \sqrt{\frac{8}{\pi}} : \sqrt{3} \approx 1 : 1.128 : 1.224}
  \end{equation}
\end{itemize}
\end{formulabox}

\begin{formulabox}[Law of Equipartition of Energy, Degrees of Freedom \& Specific Heat Ratios]
\textbf{1. Equipartition Principle:}
In thermal equilibrium at temperature $T$, the total internal energy is partitioned equally among all quadratic independent degrees of freedom ($f$), each contributing an average energy of $\mathbf{\frac{1}{2}k_B T}$ per molecule:
\begin{equation}
  \overline{E}_{\text{molecule}} = \frac{f}{2}k_B T, \quad U_{\text{molar}} = \frac{f}{2} R T \implies \mathbf{C_v = \frac{f}{2}R}, \quad \mathbf{C_p = C_v + R = \left(\frac{f}{2} + 1\right)R}, \quad \mathbf{\gamma = \frac{C_p}{C_v} = 1 + \frac{2}{f}}
\end{equation}
\textbf{2. Comprehensive Gas Classification Table:}
\begin{center}
\begin{tabularx}{\linewidth}{|l|c|c|c|c|Y|}
\hline
\textbf{Molecular Atomicity} & \textbf{DoF ($f$)} & $C_v$ & $C_p$ & $\gamma = C_p / C_v$ & \textbf{Representative Gases} \\
\hline
\textbf{Monoatomic} & $3$ trans & $\frac{3}{2}R$ & $\frac{5}{2}R$ & $\frac{5}{3} \approx 1.67$ & $\mathrm{He, Ne, Ar, Kr}$ \\
\hline
\textbf{Diatomic (Rigid, Room $T$)} & $3\text{ trans} + 2\text{ rot} = 5$ & $\frac{5}{2}R$ & $\frac{7}{2}R$ & $\frac{7}{5} = 1.40$ & $\mathrm{H_2, N_2, O_2, CO}$ \\
\hline
\textbf{Diatomic (Vibrating, High $T$)} & $5 + 2\text{ vib} = 7$ & $\frac{7}{2}R$ & $\frac{9}{2}R$ & $\frac{9}{7} \approx 1.29$ & $\mathrm{Cl_2, Br_2}$ at high $T$ \\
\hline
\textbf{Non-linear Polyatomic} & $3\text{ trans} + 3\text{ rot} = 6$ & $3R$ & $4R$ & $\frac{4}{3} \approx 1.33$ & $\mathrm{H_2O, NH_3, CH_4}$ \\
\hline
\end{tabularx}
\end{center}
\textbf{3. Non-Reacting Mixtures of Gases:}
For a mixture of $n_1$ moles of gas 1 ($C_{v1}, \gamma_1$) and $n_2$ moles of gas 2 ($C_{v2}, \gamma_2$):
\begin{equation}
  \mathbf{C_{v,\text{mix}} = \frac{n_1 C_{v1} + n_2 C_{v2}}{n_1 + n_2}}, \qquad \mathbf{C_{p,\text{mix}} = \frac{n_1 C_{p1} + n_2 C_{p2}}{n_1 + n_2}}, \qquad \mathbf{\gamma_{\text{mix}} = \frac{C_{p,\text{mix}}}{C_{v,\text{mix}}} = \frac{n_1 C_{p1} + n_2 C_{p2}}{n_1 C_{v1} + n_2 C_{v2}}}
\end{equation}
\textbf{4. Mean Free Path ($\lambda$):}
Average distance traversed by a gas molecule between successive collisions:
\begin{equation}
  \mathbf{\lambda = \frac{1}{\sqrt{2} \pi n_v d^2} = \frac{k_B T}{\sqrt{2} \pi d^2 P}}
\end{equation}
where $n_v = N/V$ is number density, $d$ is molecular diameter. Notice that $\lambda \propto T$ at constant pressure, and $\lambda \propto \frac{1}{P}$ at constant temperature.
\end{formulabox}

\section{Simple Harmonic Motion (SHM): Kinematics, Energy \& Oscillators}

\begin{conceptbox}[Mathematical Foundations of Simple Harmonic Motion]
SHM is bounded periodic motion wherein the restoring force/torque is directly proportional to displacement from the stable equilibrium position and directed towards that equilibrium:
\begin{equation}
  F(x) = -k x \implies m \frac{d^2 x}{dt^2} + k x = 0 \implies \mathbf{\frac{d^2 x}{dt^2} + \omega^2 x = 0} \quad \left(\text{where } \omega = \sqrt{\frac{k}{m}}\right)
\end{equation}
\textbf{General Solution \& Kinematic Equations:}
\begin{itemize}
  \item \textbf{Displacement:} $x(t) = A \sin(\omega t + \phi_0)$, where $A$ is amplitude, $\omega = 2\pi/T = 2\pi\nu$ is angular frequency, and $\phi_0$ is initial epoch.
  \item \textbf{Velocity:}
  \begin{equation}
    v(t) = \frac{dx}{dt} = A\omega \cos(\omega t + \phi_0) = \mathbf{\pm \omega \sqrt{A^2 - x^2}}
  \end{equation}
  $v_{\max} = \omega A$ occurs at mean position ($x=0$); $v=0$ at extreme positions ($x = \pm A$).
  \item \textbf{Acceleration:}
  \begin{equation}
    a(t) = \frac{dv}{dt} = -\omega^2 A \sin(\omega t + \phi_0) = \mathbf{-\omega^2 x}
  \end{equation}
  $a_{\max} = \omega^2 A$ occurs at extremes ($x = \pm A$); $a=0$ at mean position ($x=0$).
  \item \textbf{Phase Relationships:} Velocity leads displacement by $\pi/2$ radians ($90^\circ$); acceleration leads displacement by $\pi$ radians ($180^\circ$).
\end{itemize}
\end{conceptbox}

\begin{formulabox}[Energy Conservation in SHM: Space \& Time Averages]
\textbf{1. Instantaneous Energies as Functions of Position ($x$):}
\begin{itemize}
  \item \textbf{Potential Energy ($U$):} $U(x) = \int_0^x kx dx = \mathbf{\frac{1}{2} k x^2 = \frac{1}{2} m \omega^2 x^2}$ (assuming $U(0) = 0$).
  \item \textbf{Kinetic Energy ($K$):} $K(x) = \frac{1}{2} m v^2 = \mathbf{\frac{1}{2} m \omega^2 (A^2 - x^2)}$.
  \item \textbf{Total Mechanical Energy ($E$):}
  \begin{equation}
    \mathbf{E = K(x) + U(x) = \frac{1}{2} m \omega^2 A^2 = 2 \pi^2 m \nu^2 A^2 = \text{constant}}
  \end{equation}
  Total energy is completely independent of time $t$ and position $x$.
\end{itemize}
\textbf{2. Position where $K = U$:}
\begin{equation}
  \frac{1}{2}m\omega^2(A^2 - x^2) = \frac{1}{2}m\omega^2 x^2 \implies 2x^2 = A^2 \implies \mathbf{x = \pm \frac{A}{\sqrt{2}}}
\end{equation}
\textbf{3. Space Averages vs. Time Averages over One Cycle:}
\begin{itemize}
  \item \textbf{Time Averages ($\langle \dots \rangle_t = \frac{1}{T}\int_0^T \dots dt$):}
  \begin{equation}
    \mathbf{\langle K \rangle_t = \frac{1}{4} m \omega^2 A^2 = \frac{1}{2}E}, \qquad \mathbf{\langle U \rangle_t = \frac{1}{4} m \omega^2 A^2 = \frac{1}{2}E}
  \end{equation}
  \item \textbf{Space Averages ($\langle \dots \rangle_x = \frac{1}{2A}\int_{-A}^A \dots dx$):}
  \begin{equation}
    \mathbf{\langle K \rangle_x = \frac{2}{3}E}, \qquad \mathbf{\langle U \rangle_x = \frac{1}{3}E}
  \end{equation}
  \item \textbf{Frequency of Oscillations:} If displacement oscillates with frequency $\nu$, both kinetic energy $K(t)$ and potential energy $U(t)$ oscillate with \textbf{double frequency $2\nu$}.
\end{itemize}
\end{formulabox}

\begin{formulabox}[Mechanical Oscillators: Spring-Mass Systems \& Simple Pendulums]
\textbf{1. Spring-Block Oscillators:}
\begin{equation}
  T = 2\pi \sqrt{\frac{m}{k}}
\end{equation}
\begin{itemize}
  \item \textbf{Series Combination:} $\frac{1}{k_{\text{eq}}} = \frac{1}{k_1} + \frac{1}{k_2} \implies k_{\text{eq}} = \frac{k_1 k_2}{k_1 + k_2} \implies T = 2\pi\sqrt{\frac{m(k_1+k_2)}{k_1 k_2}}$.
  \item \textbf{Parallel Combination:} $k_{\text{eq}} = k_1 + k_2 \implies T = 2\pi\sqrt{\frac{m}{k_1 + k_2}}$.
  \item \textbf{Spring Cutting:} Spring constant is inversely proportional to length: $k \cdot L = \text{const}$. If cut into length ratio $m : n$, spring constants become $k_1 = \frac{m+n}{m}k$ and $k_2 = \frac{m+n}{n}k$.
  \item \textbf{Two-Body Oscillator Connected by Spring:} Reduced mass $\mu = \frac{m_1 m_2}{m_1 + m_2} \implies \mathbf{T = 2\pi\sqrt{\frac{\mu}{k}}}$.
\end{itemize}
\textbf{2. Simple Pendulum Derivation:}
A bob of mass $m$ suspended by a light string of length $L$. For angular displacement $\theta$:
\begin{equation}
  \tau_{\text{restoring}} = -m g L \sin\theta \approx -m g L \theta \quad (\text{for paraxial small angles } \theta \ll 1)
\end{equation}
Using $\tau = I \alpha = (m L^2) \frac{d^2\theta}{dt^2}$:
\begin{equation}
  m L^2 \frac{d^2\theta}{dt^2} + m g L \theta = 0 \implies \frac{d^2\theta}{dt^2} + \frac{g}{L}\theta = 0 \implies \mathbf{\omega = \sqrt{\frac{g}{L}}}, \quad \mathbf{T = 2\pi\sqrt{\frac{L}{g}}}
\end{equation}
\textit{Accelerating Frames \& Effective Gravity:}
\begin{itemize}
  \item Elevator accelerating upwards with $a$: $g_{\text{eff}} = g + a \implies T = 2\pi\sqrt{\frac{L}{g+a}}$ (period decreases).
  \item Elevator accelerating downwards with $a$: $g_{\text{eff}} = g - a \implies T = 2\pi\sqrt{\frac{L}{g-a}}$ (period increases; free fall $a=g \implies T \to \infty$).
  \item Horizontally accelerating vehicle: $g_{\text{eff}} = \sqrt{g^2 + a^2} \implies T = 2\pi\sqrt{\frac{L}{\sqrt{g^2 + a^2}}}$.
  \item Pendulum in liquid of density $\sigma$ (bob density $\rho$): $g_{\text{eff}} = g\left(1 - \frac{\sigma}{\rho}\right) \implies T = 2\pi\sqrt{\frac{L}{g(1 - \sigma/\rho)}}$.
\end{itemize}
\textbf{3. Seconds Pendulum:} A pendulum whose time period is exactly $T = 2.0\,\mathrm{s}$ ($1\,\mathrm{s}$ tick each half-swing):
\begin{equation}
  L = \frac{g T^2}{4\pi^2} = \frac{9.8 \times 4}{4\pi^2} \approx \mathbf{0.993\,\mathrm{m} \approx 1.0\,\mathrm{m}}
\end{equation}
\end{formulabox}

\begin{usecasebox}[Phase Circle (Phasor) Technique for Instantaneous Transit Times]
To compute the transit time between any two positions $x_1$ and $x_2$ in SHM, project the motion onto a reference circle of radius $A$ rotating counterclockwise at angular velocity $\omega$:
\begin{equation}
  x = A \sin\theta, \quad \theta = \omega t \implies t = \frac{\Delta\theta}{\omega} = \frac{\Delta\theta}{2\pi} T
\end{equation}
\textit{High-Yield Benchmarks (Starting from Mean Position $x=0$):}
\begin{itemize}
  \item From $x = 0$ to $x = \frac{A}{2}$: $\theta_1 = 0 \to \theta_2 = 30^\circ = \frac{\pi}{6} \implies \mathbf{t = \frac{T}{12}}$.
  \item From $x = \frac{A}{2}$ to $x = A$: $\Delta\theta = 90^\circ - 30^\circ = 60^\circ = \frac{\pi}{3} \implies \mathbf{t = \frac{T}{6}}$.
  \item Notice: Traveling the outer half of amplitude takes \textbf{twice as long} as the inner half ($T/6$ vs $T/12$) due to deceleration.
  \item From $x = 0$ to $x = \frac{A}{\sqrt{2}}$: $\theta = 45^\circ = \frac{\pi}{4} \implies \mathbf{t = \frac{T}{8}}$.
\end{itemize}
\end{usecasebox}

\section{Wave Motion, Stretched Strings \& Organ Pipes}

\begin{conceptbox}[Kinematics of 1D Harmonic Progressive Waves]
A 1D transverse or longitudinal sinusoidal traveling wave propagating along the $+x$ direction:
\begin{equation}
  \mathbf{y(x,t) = A \sin(k x - \omega t + \phi) = A \sin\left[2\pi\left(\frac{x}{\lambda} - \nu t\right) + \phi\right]}
\end{equation}
where $k = \frac{2\pi}{\lambda}$ is the angular wave number, $\omega = 2\pi\nu$ is angular frequency.
\begin{itemize}
  \item \textbf{Wave Speed ($v$):} Speed of propagation of constant phase profile:
  \begin{equation}
    k x - \omega t = \text{const} \implies k \frac{dx}{dt} - \omega = 0 \implies \mathbf{v = \frac{\omega}{k} = \nu \lambda}
  \end{equation}
  \item \textbf{Particle Velocity ($v_p$) vs. Wave Velocity ($v$):}
  \begin{equation}
    v_p = \frac{\partial y}{\partial t} = -\omega A \cos(k x - \omega t), \qquad \frac{\partial y}{\partial x} = k A \cos(k x - \omega t)
  \end{equation}
  \begin{equation}
    \mathbf{v_p = -v \left(\frac{\partial y}{\partial x}\right) = -v \times (\text{Wave profile slope})}
  \end{equation}
  If wave moves right ($v>0$) and slope is positive ($\partial y/\partial x > 0$), particle moves downwards ($v_p < 0$).
  \item \textbf{Particle Acceleration:} $a_p = \frac{\partial^2 y}{\partial t^2} = -\omega^2 y$.
\end{itemize}
\end{conceptbox}

\begin{formulabox}[Speeds of Transverse Waves on Strings \& Sound in Gases]
\textbf{1. Transverse Wave on Stretched String:}
\begin{equation}
  \mathbf{v = \sqrt{\frac{T}{\mu}}}
\end{equation}
where $T$ is tension ($\mathrm{N}$), and $\mu = \frac{M}{L} = \rho A$ is linear mass density ($\mathrm{kg/m}$).
\textit{Derivation via Centripetal Equilibrium of Arc Element:}
Consider a small curved element of string of length $\Delta l = R \Delta\theta$ moving with wave speed $v$ along a circular crest of radius $R$. Mass of element $\Delta m = \mu \Delta l = \mu R \Delta\theta$.
The inward restoring component of tension:
\begin{equation}
  F_{\text{inward}} = 2 T \sin\left(\frac{\Delta\theta}{2}\right) \approx T \Delta\theta
\end{equation}
Equating to centripetal force:
\begin{equation}
  T \Delta\theta = \frac{\Delta m \cdot v^2}{R} = \frac{(\mu R \Delta\theta) v^2}{R} = \mu v^2 \Delta\theta \implies \mathbf{v = \sqrt{\frac{T}{\mu}}}
\end{equation}
\textbf{2. Speed of Longitudinal Waves (Sound) in Fluids \& Gases:}
\begin{itemize}
  \item General relation: $v = \sqrt{\frac{B}{\rho}}$, where $B$ is bulk modulus of elasticity.
  \item \textbf{Newton's Isothermal Formula (Erroneous):} Newton assumed sound compressions/rarefactions occur isothermally ($B = P$).
  \begin{equation}
    v_{\text{Newton}} = \sqrt{\frac{P}{\rho}} \approx \sqrt{\frac{1.013 \times 10^5}{1.293}} \approx 280\,\mathrm{m/s} \quad (\approx 15\% \text{ below experimental } 332\,\mathrm{m/s})
  \end{equation}
  \item \textbf{Laplace's Adiabatic Correction (Accurate):} Compressions and rarefactions occur so rapidly that heat exchange is zero ($dQ=0 \implies B_{\text{adiabatic}} = \gamma P$):
  \begin{equation}
    \mathbf{v = \sqrt{\frac{\gamma P}{\rho}} = \sqrt{\frac{\gamma R T}{M_{\text{molar}}}}}
  \end{equation}
  For air ($\gamma \approx 1.4, M \approx 28.8 \times 10^{-3}\,\mathrm{kg/mol}$ at $273\,\mathrm{K}$): $v = 332\,\mathrm{m/s}$ (perfect agreement).
  \item \textbf{Dependence on Environmental Factors:}
  \begin{itemize}
    \item \textit{Pressure:} Invariant of pressure at constant temperature (since $P/\rho = RT/M = \text{const}$).
    \item \textit{Temperature:} $v \propto \sqrt{T} \implies v_t \approx v_0 + 0.61 t$ (for small temperatures $t^\circ\mathrm{C}$).
    \item \textit{Humidity:} Moist air is less dense than dry air ($\rho_{\text{moist}} < \rho_{\text{dry}}$), so speed of sound \textbf{increases} with humidity.
  \end{itemize}
\end{itemize}
\end{formulabox}

\begin{formulabox}[Standing Waves: Strings \& Organ Pipes with End Corrections]
\textbf{1. Superposition \& Standing Wave Equation on String Fixed at Both Ends:}
Two identical waves traveling in opposite directions: $y_1 = A\sin(kx - \omega t)$ and $y_2 = A\sin(kx + \omega t)$.
\begin{equation}
  y_{\text{net}} = y_1 + y_2 = \mathbf{[2A \sin(kx)] \cos(\omega t)}
\end{equation}
\begin{itemize}
  \item \textbf{Nodes ($y=0$):} $\sin(kx) = 0 \implies kx = n\pi \implies \mathbf{x = \frac{n\lambda}{2}}$ ($n = 0, 1, 2, \dots$).
  \item \textbf{Antinodes ($|y|=2A$):} $|\sin(kx)| = 1 \implies \mathbf{x = (2n-1)\frac{\lambda}{4}}$.
  \item Boundary condition: nodes at $x = 0$ and $x = L \implies L = n\frac{\lambda}{2} \implies \mathbf{\nu_n = n \left(\frac{v}{2L}\right) = n \nu_1}$ (\textbf{All harmonics present}: $1 : 2 : 3 : \dots$).
\end{itemize}
\textbf{2. Closed Organ Pipe (One End Closed, One End Open):}
Boundary condition: Node at closed end, Antinode at open end. Length $L = (2n-1)\frac{\lambda}{4}$:
\begin{equation}
  \mathbf{\nu_n = (2n-1)\frac{v}{4L} = (2n-1)\nu_1} \qquad \text{where fundamental } \nu_1 = \frac{v}{4L}
\end{equation}
\textbf{Only odd harmonics exist} in a closed organ pipe ($\nu_1 : \nu_3 : \nu_5 = 1 : 3 : 5 : \dots$). No even harmonics.
\textbf{3. Open Organ Pipe (Both Ends Open):}
Boundary condition: Antinodes at both open ends. Length $L = n\frac{\lambda}{2}$:
\begin{equation}
  \mathbf{\nu_n = n \frac{v}{2L} = n \nu_1} \qquad \text{where fundamental } \nu_1 = \frac{v}{2L}
\end{equation}
\textbf{All harmonics exist} ($\nu_1 : \nu_2 : \nu_3 = 1 : 2 : 3 : \dots$). Tone is richer/sweeter than closed pipe.
\textbf{4. End Correction ($e$):}
Antinode forms slightly outside the physical open tube boundary due to finite acoustic inertia of moving air:
\begin{equation}
  e \approx 0.6 r \quad (\text{where } r \text{ is tube inner radius})
\end{equation}
\begin{itemize}
  \item For Closed Pipe: $L_{\text{eff}} = L + e = L + 0.6r$.
  \item For Open Pipe: $L_{\text{eff}} = L + 2e = L + 1.2r$.
  \item \textbf{Resonance Tube Method to eliminate $e$:}
  Using 1st resonance length $l_1$ and 2nd resonance length $l_2$ for a tuning fork of frequency $\nu$:
  \begin{equation}
    l_1 + e = \frac{\lambda}{4}, \qquad l_2 + e = \frac{3\lambda}{4} \implies l_2 - l_1 = \frac{\lambda}{2} \implies \mathbf{\lambda = 2(l_2 - l_1)}
  \end{equation}
  \begin{equation}
    \mathbf{v = 2\nu(l_2 - l_1)}, \qquad \mathbf{e = \frac{l_2 - 3l_1}{2}}
  \end{equation}
\end{itemize}
\end{formulabox}

\begin{formulabox}[Acoustic Beats: Analytical Derivation]
Consider superposition of two harmonic sound waves of equal amplitude $A$ and slightly differing frequencies $\nu_1, \nu_2$ ($\omega_1 \approx \omega_2$):
\begin{equation}
  y_1 = A\sin(\omega_1 t), \qquad y_2 = A\sin(\omega_2 t)
\end{equation}
\begin{equation}
  y_{\text{net}} = y_1 + y_2 = 2A \cos\left(\frac{\omega_1 - \omega_2}{2}t\right) \sin\left(\frac{\omega_1 + \omega_2}{2}t\right) = A_{\text{mod}}(t) \sin(\omega_{\text{avg}} t)
\end{equation}
The modulated amplitude is $A_{\text{mod}}(t) = 2A \cos(\pi(\nu_1 - \nu_2)t)$.
Intensity $I \propto A_{\text{mod}}^2 \propto \cos^2(\pi(\nu_1 - \nu_2)t) = \frac{1 + \cos(2\pi(\nu_1 - \nu_2)t)}{2}$.
Maximum intensity occurs whenever $2\pi(\nu_1 - \nu_2)t = n\pi \implies t = \frac{n}{|\nu_1 - \nu_2|}$.
Time interval between consecutive waxing/waning: $\Delta t = \frac{1}{|\nu_1 - \nu_2|}$.
\begin{equation}
  \mathbf{\nu_{\text{beat}} = |\nu_1 - \nu_2|}
\end{equation}
\textit{Tuning Fork Rules:}
\begin{itemize}
  \item Loading fork with wax increases inertia $\implies \nu$ decreases.
  \item Filing fork prongs decreases mass $\implies \nu$ increases.
\end{itemize}
\end{formulabox}

\begin{trapbox}[Syllabus Boundary Alert: Complete Deletion of Doppler Effect]
\textbf{NMC REVISION WARNING FOR NEET UG:}
Under the revised NMC curriculum for NEET UG 2027:
\begin{itemize}
  \item \textbf{STRICTLY DELETED:} The \textbf{Doppler Effect in Sound Waves} (apparent frequency shifts $f' = f\frac{v \pm v_0}{v \mp v_s}$, moving observers, moving sources, echo reflection from cliffs, acoustic shockwaves) has been \textbf{completely eliminated} from the syllabus.
  \item \textbf{EXAM ORIENTATION:} Candidates must not solve Doppler frequency formulas. Instead, master Standing Waves, Harmonics in Organ Pipes, Resonance Tube end corrections, and Beats.
\end{itemize}
\end{trapbox}

\section{Master High-Yield Problem Archetypes \& Worked Solutions}

\begin{usecasebox}[Archetype 1: Ice-Water-Steam Calorimetry Mixing Equilibrium]
\textbf{Problem:} $10\,\mathrm{g}$ of ice at $-10^\circ\mathrm{C}$ is mixed with $10\,\mathrm{g}$ of steam at $100^\circ\mathrm{C}$ in an insulated calorimeter of negligible heat capacity. Determine the final equilibrium temperature and the constitution of the mixture.
\textbf{Step-by-Step Solution:}
\begin{enumerate}
  \item Heat required to bring ice from $-10^\circ\mathrm{C}$ to liquid water at $100^\circ\mathrm{C}$:
  \begin{itemize}
    \item Warm ice to $0^\circ\mathrm{C}$: $Q_1 = m_{\text{ice}} s_{\text{ice}} \Delta T = 10 \times 0.5 \times 10 = 50\,\mathrm{cal}$.
    \item Melt ice at $0^\circ\mathrm{C}$: $Q_2 = m_{\text{ice}} L_f = 10 \times 80 = 800\,\mathrm{cal}$.
    \item Warm water to $100^\circ\mathrm{C}$: $Q_3 = m_{\text{ice}} s_w \Delta T = 10 \times 1.0 \times 100 = 1000\,\mathrm{cal}$.
    \item Total heat absorbed to reach $100^\circ\mathrm{C}$ water: $Q_{\text{absorb}} = 50 + 800 + 1000 = \mathbf{1850\,\mathrm{cal}}$.
  \end{itemize}
  \item Total latent heat available from steam condensing at $100^\circ\mathrm{C}$:
  \begin{equation}
    Q_{\text{condense, total}} = m_{\text{steam}} L_v = 10 \times 540 = \mathbf{5400\,\mathrm{cal}}
  \end{equation}
  \item Notice that $Q_{\text{condense, total}} = 5400\,\mathrm{cal} > Q_{\text{absorb}} = 1850\,\mathrm{cal}$.
  Therefore, not all steam condenses! The mixture reaches dynamic equilibrium at \textbf{$T_{\text{eq}} = 100^\circ\mathrm{C}$}.
  \item Mass of steam condensed to supply $1850\,\mathrm{cal}$:
  \begin{equation}
    m_{\text{condensed}} = \frac{1850\,\mathrm{cal}}{540\,\mathrm{cal/g}} \approx 3.43\,\mathrm{g}
  \end{equation}
  \item Surviving steam: $m_{\text{steam, remaining}} = 10 - 3.43 = \mathbf{6.57\,\mathrm{g}}$.
  \item Total water at $100^\circ\mathrm{C}$: $m_{\text{water}} = 10 + 3.43 = \mathbf{13.43\,\mathrm{g}}$.
\end{enumerate}
\end{usecasebox}

\begin{usecasebox}[Archetype 2: Mixture of Distinct Polyatomic \& Monoatomic Gases]
\textbf{Problem:} $2$ moles of Helium (monoatomic, $\gamma_1 = 5/3$) are mixed with $3$ moles of Oxygen (diatomic, rigid, $\gamma_2 = 7/5$). Calculate $C_{v,\text{mix}}$, $C_{p,\text{mix}}$, and the adiabatic exponent $\gamma_{\text{mix}}$.
\textbf{Step-by-Step Solution:}
\begin{enumerate}
  \item Molar heat capacities of individual gases:
  \begin{itemize}
    \item Helium (monoatomic): $C_{v1} = \frac{3}{2}R$, $C_{p1} = \frac{5}{2}R$.
    \item Oxygen (diatomic): $C_{v2} = \frac{5}{2}R$, $C_{p2} = \frac{7}{2}R$.
  \end{itemize}
  \item Total moles: $n_{\text{total}} = 2 + 3 = 5\,\mathrm{mol}$.
  \item Equivalent $C_{v,\text{mix}}$:
  \begin{equation}
    C_{v,\text{mix}} = \frac{n_1 C_{v1} + n_2 C_{v2}}{n_1 + n_2} = \frac{2 \left(\frac{3}{2}R\right) + 3 \left(\frac{5}{2}R\right)}{5} = \frac{3R + 7.5R}{5} = \frac{10.5R}{5} = \mathbf{2.1 R} = \frac{21}{10}R
  \end{equation}
  \item Equivalent $C_{p,\text{mix}}$:
  \begin{equation}
    C_{p,\text{mix}} = C_{v,\text{mix}} + R = 2.1 R + R = \mathbf{3.1 R} = \frac{31}{10}R
  \end{equation}
  \item Equivalent adiabatic index $\gamma_{\text{mix}}$:
  \begin{equation}
    \gamma_{\text{mix}} = \frac{C_{p,\text{mix}}}{C_{v,\text{mix}}} = \frac{3.1 R}{2.1 R} = \frac{31}{21} \approx \mathbf{1.476}
  \end{equation}
\end{enumerate}
\end{usecasebox}

\begin{usecasebox}[Archetype 3: Resonance Tube End Correction \& Speed of Sound]
\textbf{Problem:} A tuning fork of frequency $\nu = 512\,\mathrm{Hz}$ excites the first two resonances in a resonance water tube at lengths $l_1 = 15.2\,\mathrm{cm}$ and $l_2 = 48.0\,\mathrm{cm}$. Compute the speed of sound in air and the inner pipe radius $r$.
\textbf{Step-by-Step Solution:}
\begin{enumerate}
  \item The fundamental difference between resonances isolates the half-wavelength:
  \begin{equation}
    \frac{\lambda}{2} = l_2 - l_1 = 48.0\,\mathrm{cm} - 15.2\,\mathrm{cm} = 32.8\,\mathrm{cm} \implies \lambda = 65.6\,\mathrm{cm} = 0.656\,\mathrm{m}
  \end{equation}
  \item Speed of sound in air:
  \begin{equation}
    v = \nu \lambda = 512\,\mathrm{s}^{-1} \times 0.656\,\mathrm{m} = \mathbf{335.87\,\mathrm{m/s}}
  \end{equation}
  \item End correction $e$:
  \begin{equation}
    e = \frac{l_2 - 3l_1}{2} = \frac{48.0 - 3(15.2)}{2} = \frac{48.0 - 45.6}{2} = \frac{2.4}{2} = \mathbf{1.2\,\mathrm{cm}}
  \end{equation}
  \item Inner radius $r$ of the tube using $e \approx 0.6 r$:
  \begin{equation}
    r = \frac{e}{0.6} = \frac{1.2\,\mathrm{cm}}{0.6} = \mathbf{2.0\,\mathrm{cm}}
  \end{equation}
\end{enumerate}
\end{usecasebox}

\end{document}
